\chapter{Design Patterns and Anti-Patterns in FRFP Systems}
\label{ch:frfp-patterns}

This chapter distils the case studies of Volume~II into a small set of
\emph{design patterns} and \emph{anti-patterns} for FRFP systems.
The goal is a reusable pattern language:
a way to recognise structural moves that tend to go well or badly across
domains, and to apply those moves when designing new systems.

We write patterns in the same vocabulary as Volume~I:

\begin{itemize}
  \item explicit space \(E\) and explicit pipelines in \(N \ltimes E\);
  \item tacit space \(T\) and tacit states of humans and institutions;
  \item boundaries where explicit artefacts cross into tacit endorsement;
  \item safe horizons that bound how far runs can proceed without
        re-grounding;
  \item institutional operators that shape which runs are admissible.
\end{itemize}

Each pattern is illustrated (implicitly or explicitly) by one or more
case studies in this volume:
AlphaFold, DR screening AI, Waymo, Watson, Zillow, COMPAS, Project Vend,
Mata v.\ Avianca, the London Whale, and the platform-scale systems.

\section{Pattern P1: Narrow, Strong Explicit Semantics}
\label{pattern:p1}

\textbf{Name.} Narrow, Strong Explicit Semantics.

\textbf{Intent.} Restrict the AI system to tasks where the mapping it
implements in explicit space is well-defined, externally evaluable, and
closely aligned with the tacit notion of correctness for that task.

\subsection*{Structure}

\begin{itemize}
  \item Choose a task with a clear, domain-accepted definition:
        \begin{itemize}
          \item e.g.\ ``predict protein 3D structure from sequence'',
                ``detect more-than-mild diabetic retinopathy from a
                fundus image'', ``classify a stop sign''.
        \end{itemize}
  \item Ensure ground truth is available at scale:
        experimental structures, specialist-graded images, labelled
        scenarios, etc.
  \item Define explicit metrics that are good proxies for tacit
        correctness:
        RMSD, sensitivity/specificity, error rates on held-out data.
  \item Train and evaluate models \emph{only} within this well-defined
        semantic scope.
\end{itemize}

\subsection*{Examples}

\begin{itemize}
  \item AlphaFold:
        sequence $\to$ static structure, evaluated in CASP.
  \item DR screening AI:
        fundus image $\to$ mtmDR/no-mtmDR/ungradable, evaluated
        against masked specialist grading.
  \item Core perception modules in Waymo:
        sensor data $\to$ object detections and lane geometry.
\end{itemize}

\subsection*{Consequences}

\begin{itemize}
  \item Explicit metrics track the domain’s tacit notion of correctness.
  \item Failures are mainly statistical (noise, distribution shift),
        not semantic (wrong task).
  \item It is easier to reason formally about safe horizons, because
        the task semantics do not mutate as the system scales.
\end{itemize}

\subsection*{When to use}

\begin{itemize}
  \item In high-stakes domains where a narrow, well-defined subtask is
        available.
  \item As the foundation on top of which more complex, tacit-heavy
        workflows are built.
\end{itemize}

\section{Pattern P2: Explicit Uncertainty and ``I Don't Know''}
\label{pattern:p2}

\textbf{Name.} Explicit Uncertainty and ``I Don't Know''.

\textbf{Intent.} Design explicit pipelines so that they can express when
they are unreliable, and route those cases back into tacit space instead
of forcing a brittle decision.

\subsection*{Structure}

\begin{itemize}
  \item Augment outputs with explicit uncertainty or quality measures:
        \begin{itemize}
          \item per-residue confidence scores (pLDDT, PAE);
          \item image-quality flags and ``ungradable'' outcomes;
          \item confidence intervals, anomaly flags, or out-of-domain
                detection.
        \end{itemize}
  \item Treat ``I don't know'' or ``ungradable'' as first-class outputs,
        \emph{not} as errors.
  \item Define routing rules:
        when uncertainty exceeds a threshold, control passes to human
        tacit judgement or to a different, more conservative pipeline.
\end{itemize}

\subsection*{Examples}

\begin{itemize}
  \item AlphaFold’s confidence scores and PAE maps:
        users interpret low-confidence regions differently from
        high-confidence cores.
  \item DR screening AI’s ``ungradable'' outcome:
        poor images trigger referral rather than a false negative.
  \item Waymo’s safe-stop behaviours when perception is degraded or
        conditions leave the ODD.
\end{itemize}

\subsection*{Consequences}

\begin{itemize}
  \item Reduces catastrophic errors in regions where explicit semantics
        are fragile.
  \item Forces boundary events back into \(T\) when \(E\) loses grip.
  \item Makes it easier for institutional operators to specify when AI
        may or may not be used autonomously.
\end{itemize}

\subsection*{When to use}

\begin{itemize}
  \item Whenever models will unavoidably face out-of-distribution cases
        or ambiguous inputs.
  \item In any domain where false confidence is more dangerous than
        equivocation (law, medicine, finance, safety-critical control).
\end{itemize}

\section{Pattern P3: Boundary-Centred Interfaces}
\label{pattern:p3}

\textbf{Name.} Boundary-Centred Interfaces.

\textbf{Intent.} Design systems so that boundaries---points where humans
endorse or commit---are explicit, instrumented, and supported by the
right information, instead of being implicit in a UI or process.

\subsection*{Structure}

\begin{itemize}
  \item Identify boundary events:
        filings, orders, diagnoses, trade approvals, deployments,
        policy adoptions.
  \item For each boundary, specify:
        \begin{itemize}
          \item who is authorised to endorse;
          \item what artefacts they see (raw data, model outputs,
                uncertainty, alternatives);
          \item what gets logged (identity, time, context, rationale).
        \end{itemize}
  \item Design UIs and workflows around boundaries:
        the interface foregrounds the boundary decision, not just model
        outputs.
  \item Treat signatures and approvals as semantic operations:
        transitions from ``proposal'' to ``endorsed'', visible in the
        system state.
\end{itemize}

\subsection*{Examples}

\begin{itemize}
  \item DR screening workflows:
        AI result $\to$ documented screening status / referral.
  \item Waymo:
        system-level decisions about where and when to operate
        driverlessly, with explicit ODD declarations and permits.
  \item The legal case-study counterfactuals:
        AI-generated drafts, but human-verified citations at filing.
\end{itemize}

\subsection*{Consequences}

\begin{itemize}
  \item Makes it harder for AI systems to silently capture correctness.
  \item Creates a clear audit trail of tacit endorsements.
  \item Supports institutional oversight: committees can review
        boundary events, not just raw logs.
\end{itemize}

\subsection*{When to use}

\begin{itemize}
  \item For any system affecting rights, money, safety, or scientific
        claims.
  \item Wherever existing professional norms already treat certain
        actions (signing, prescribing, executing a trade) as serious
        commitments.
\end{itemize}

\section{Pattern P4: Engineered Safe Horizons}
\label{pattern:p4}

\textbf{Name.} Engineered Safe Horizons.

\textbf{Intent.} Make the limits of a system’s operation in time, scale,
and domain explicit and enforce them with mechanisms that force
re-grounding in tacit judgement.

\subsection*{Structure}

\begin{itemize}
  \item Define horizons along three axes:
        \begin{itemize}
          \item temporal (how long between thorough reviews);
          \item scale (how many cases, trades, patients, rides, or
                items affected);
          \item domain (how far from the context where the system was
                validated).
        \end{itemize}
  \item Associate each horizon with triggers:
        caps, trial endpoints, or conditions under which operation must
        pause or shift mode.
  \item Define re-grounding procedures:
        audits, revalidation, redesign, or human-only periods.
  \item Use phase-gated deployment:
        pilot $\to$ limited expansion $\to$ broader roll-out, with
        explicit go/no-go decisions between phases.
\end{itemize}

\subsection*{Examples}

\begin{itemize}
  \item DR screening AI:
        yearly re-screening intervals; post-market surveillance; limits
        on indications and camera types.
  \item Waymo:
        operational design domains (cities, weather, road types);
        gradual expansion; service suspension in out-of-ODD conditions.
  \item Scientific workflows with AlphaFold:
        predictions feed into experiments; no deployment of predictions
        as clinical facts without separate validation.
\end{itemize}

\subsection*{Consequences}

\begin{itemize}
  \item Reduces the risk of slow, unnoticed drift.
  \item Forces institutions to periodically reflect and adjust.
  \item Provides natural points to change objectives, models, and
        governance in light of new evidence.
\end{itemize}

\subsection*{When to use}

\begin{itemize}
  \item For any long-horizon, high-impact system (clinical programmes,
        trading desks, platform-wide algorithms).
  \item Especially when scaling pilots to full production.
\end{itemize}

\section{Pattern P5: Governance as a First-Class Operator}
\label{pattern:p5}

\textbf{Name.} Governance as a First-Class Operator.

\textbf{Intent.} Treat governance as a semantic layer with its own
operators on tacit states and admissible runs, not as an afterthought
described only in prose or policy PDF files.

\subsection*{Structure}

\begin{itemize}
  \item Identify governance bodies:
        internal committees, external regulators, boards, professional
        societies, ethics panels.
  \item For each, specify:
        \begin{itemize}
          \item what information they see (metrics, incidents,
                case studies, whistleblower reports);
          \item what operations they can perform:
                approve, suspend, restrict, mandate changes, redefine
                objectives;
          \item how frequently they act and under what triggers.
        \end{itemize}
  \item Ensure governance is structurally distinct from the teams that
        build and operate the system, even if they share information.
  \item Model governance decisions as changes in:
        \begin{itemize}
          \item boundaries (what AI can/cannot decide);
          \item horizons (caps, trial lengths, ODDs);
          \item acceptable explicit tasks and objectives.
        \end{itemize}
\end{itemize}

\subsection*{Examples}

\begin{itemize}
  \item FDA and similar regulators for DR screening AI:
        explicit indications-for-use, performance targets, and
        post-market requirements.
  \item EMBL--EBI and the structural biology community around AF-DB:
        shared norms on how AlphaFold predictions should be used.
  \item City and state regulators for Waymo:
        permits, geographic constraints, and operational restrictions.
\end{itemize}

\subsection*{Consequences}

\begin{itemize}
  \item Makes institutional non-automation visible and analysable.
  \item Gives a formal place to value judgements and political choices.
  \item Allows comparison of governance designs across domains.
\end{itemize}

\subsection*{When to use}

\begin{itemize}
  \item Always, for systems that affect large populations or high-stakes
        domains.
  \item Even for internal tools, where a ``governance board'' might be
        a cross-functional risk committee rather than a regulator.
\end{itemize}

\bigskip

We now contrast these positive patterns with recurrent anti-patterns from
the failure case studies.

\section{Anti-Pattern A1: Advisory-in-Name-Only}
\label{antipattern:a1}

\textbf{Name.} Advisory-in-Name-Only.

\textbf{Smell.} The system is officially ``decision support'', but in
practice its outputs are rarely overruled and often determine outcomes.

\subsection*{Structure}

\begin{itemize}
  \item Model outputs are integrated into workflows as default
        decisions:
        risk scores that drive bail decisions, hiring scores that
        determine interview lists, recommendations that become offers.
  \item Human ``overrides'' are possible but rare, costly, or
        discouraged.
  \item Interfaces emphasise model scores, not uncertainty or
        alternative reasoning.
\end{itemize}

\subsection*{Consequences}

\begin{itemize}
  \item De facto boundary capture:
        correctness migrates from tacit space \(T\) to the explicit
        pipeline \(N \ltimes E\), contradicting stated governance.
  \item Formal accountability remains with humans, but their practical
        control is diminished.
  \item Failures are attributed to ``human error'' despite being driven
        structurally by model design and integration.
\end{itemize}

\subsection*{Examples}

\begin{itemize}
  \item COMPAS scores in courts which, while officially advisory, often
        significantly influence bail and sentencing.
  \item Algorithmic ranking in hiring or admissions that effectively
        determines the candidate pool.
\end{itemize}

\subsection*{Refactoring toward patterns}

\begin{itemize}
  \item Introduce boundary-centred interfaces:
        require structured justification for following or overriding
        the tool, logged for review.
  \item Narrow the explicit task and ensure strong explicit semantics
        (Pattern~P1) and uncertainty signalling (Pattern~P2).
\end{itemize}

\section{Anti-Pattern A2: Metric Worship and Explicit-Only Governance}
\label{antipattern:a2}

\textbf{Name.} Metric Worship and Explicit-Only Governance.

\textbf{Smell.} Decision-makers treat dashboards and metrics as
authoritative, with little attention to their semantics, limits, or
unmodelled effects.

\subsection*{Structure}

\begin{itemize}
  \item A small set of key metrics (VaR, engagement, accuracy, NPS,
        ROC-AUC) dominates governance conversations.
  \item Decision-makers focus on whether metrics are within target
        bands, not on whether the metrics still measure what matters.
  \item Qualitative and tacit signals (frontline feedback, user
        discomfort, external concerns) are dismissed as ``anecdotal''
        or noisy.
\end{itemize}

\subsection*{Consequences}

\begin{itemize}
  \item Model mis-specification and distribution shift go unnoticed
        until large failures occur.
  \item Governance is effectively automated: institutional operators
        become operators on metrics rather than on real-world states.
  \item Incentives skew toward metric gaming and cosmetic fixes.
\end{itemize}

\subsection*{Examples}

\begin{itemize}
  \item Zillow Offers scaling based on backtested pricing metrics while
        ignoring selection effects and regime shifts.
  \item London Whale relying on a flawed VaR model that understated
        risk but was accepted because the numbers looked good.
  \item Watson Health emphasising demo performance and marketing
        narratives over prospective clinical evidence.
\end{itemize}

\subsection*{Refactoring toward patterns}

\begin{itemize}
  \item Make governance a first-class operator (Pattern~P5) with
        independent tacit input.
  \item Use metrics as inputs, not replacements, for boundary decisions.
  \item Introduce safe horizons (Pattern~P4) tied to independent
        evaluations, not just internal metrics.
\end{itemize}

\section{Anti-Pattern A3: One-Shot Validation, Infinite Deployment}
\label{antipattern:a3}

\textbf{Name.} One-Shot Validation, Infinite Deployment.

\textbf{Smell.} The system passes an initial validation (pilot, test,
audit) and then is deployed at scale indefinitely, with no explicit
plan for revalidation or sunset.

\subsection*{Structure}

\begin{itemize}
  \item Initial testing on a fixed dataset or limited pilot population.
  \item A one-time ``go'' decision leads to wide or permanent
        deployment.
  \item Monitoring reduces to passive metrics; there are no hard
        endpoints or caps that \emph{force} a strategic review.
\end{itemize}

\subsection*{Consequences}

\begin{itemize}
  \item The system gradually drifts away from the conditions under
        which it was validated.
  \item Accumulated effects (capital exposure, reliance, behaviour
        change) are not revisited until visible failure.
  \item It becomes politically and economically difficult to pause or
        decommission the system, even when evidence of misalignment
        appears.
\end{itemize}

\subsection*{Examples}

\begin{itemize}
  \item Zillow Offers scaling rapidly after initial model success,
        without clear phase gates tied to external conditions.
  \item Long-lived use of risk scores or heuristic tools in law,
        finance, or medicine, long after their original evidence base
        has aged or the context has shifted.
\end{itemize}

\subsection*{Refactoring toward patterns}

\begin{itemize}
  \item Replace one-shot validation with engineered safe horizons
        (Pattern~P4): pilots with explicit endpoints, caps, and
        revalidation rules.
  \item Tie horizon extensions to independent evidence, not just
        internal enthusiasm.
\end{itemize}

\section{Anti-Pattern A4: Workflow-Led Design Without Semantics}
\label{antipattern:a4}

\textbf{Name.} Workflow-Led Design Without Semantics.

\textbf{Smell.} The system is specified in terms of phases, swimlanes,
and diagrams, but there is no clear account of what the AI components
in fact compute, or what states and boundaries mean.

\subsection*{Structure}

\begin{itemize}
  \item Design documents emphasise:
        \begin{itemize}
          \item who talks to whom;
          \item which tickets are created;
          \item which committees review what, when.
        \end{itemize}
  \item There is little or no:
        \begin{itemize}
          \item formal description of explicit tasks and mappings;
          \item explicit modelling of tacit states;
          \item representation of boundaries as state transitions.
        \end{itemize}
\end{itemize}

\subsection*{Consequences}

\begin{itemize}
  \item It is difficult to reason about correctness or failure modes:
        everything is in narrative form.
  \item Workflows become brittle when context changes; there is no
        semantic core to adapt.
  \item Governance devolves into checking boxes rather than evaluating
        runs in \(N \ltimes E\).
\end{itemize}

\subsection*{Examples}

\begin{itemize}
  \item The initial V01 framework for long-horizon collaboration (the
        meta-case study):
        rich in workflow concepts, poor in semantic structure.
  \item Many contemporary ``AI governance'' and ``responsible AI''
        documents that specify processes but not semantics.
\end{itemize}

\subsection*{Refactoring toward patterns}

\begin{itemize}
  \item Start from explicit and tacit spaces (Steps~2 and~3 in
        Chapter~\ref{ch:frfp-method}), not from swimlanes.
  \item Introduce explicit semantics (Pattern~P1) and boundaries
        (Pattern~P3) before designing workflows.
  \item Treat governance as a semantic layer (Pattern~P5), not just a
        flowchart box.
\end{itemize}

\section{Anti-Pattern A5: Tacit Erosion by Over-Automation}
\label{antipattern:a5}

\textbf{Name.} Tacit Erosion by Over-Automation.

\textbf{Smell.} Over time, humans become less able or willing to
exercise independent judgement, because AI systems handle more of the
cognitive work.

\subsection*{Structure}

\begin{itemize}
  \item AI tools provide increasingly complete drafts, diagnoses,
        recommendations, or decisions.
  \item Human users are nominally ``in the loop'' but:
        \begin{itemize}
          \item rarely override the system;
          \item lose practice in underlying skills;
          \item face time pressure that incentivises acceptance.
        \end{itemize}
  \item Training and evaluation assume stable human capabilities, even
        as those capabilities change under AI influence.
\end{itemize}

\subsection*{Consequences}

\begin{itemize}
  \item When AI systems fail or encounter novel cases, human backstops
        are weaker than assumed.
  \item Organisational claims that ``humans are ultimately
        responsible'' become less realistic over time.
  \item Tacit norms (e.g.\ about verification, scepticism) decay.
\end{itemize}

\subsection*{Examples}

\begin{itemize}
  \item Lawyers relying on LLMs for research and drafting, leading to
        hallucinated citations and weakened verification norms
        (\emph{Mata v.\ Avianca} and follow-ons).
  \item Potential long-run effects of heavy LLM usage on students'
        writing, critical reading, or coding skills.
\end{itemize}

\subsection*{Refactoring toward patterns}

\begin{itemize}
  \item Design boundaries and interfaces (Pattern~P3) that require
        substantive human engagement at key points.
  \item Use safe horizons (Pattern~P4) to limit how far AI-driven
        workflows can proceed without deeper human involvement.
  \item Explicitly model tacit skills and norms as part of \(T\) and
        treat their preservation as a design objective.
\end{itemize}

\section{Using the Pattern Language}

This chapter offers a compact language for describing FRFP systems in
practice:

\begin{itemize}
  \item Patterns P1--P5 describe structural moves that appear again and
        again in successful or robust systems.
  \item Anti-patterns A1--A5 capture recurrent failure modes across
        very different domains.
\end{itemize}

When analysing or designing a system, you can use this language in two
ways:

\begin{enumerate}
  \item \textbf{Diagnostic.}
        Walk through the anti-patterns and ask:
        \begin{itemize}
          \item Do we claim to be ``advisory'' while behaving
                otherwise (A1)?
          \item Are we governing primarily by metrics (A2)?
          \item Are we treating validation as a one-time gate (A3)?
          \item Are our documents workflow-led rather than semantic
                (A4)?
          \item Are we assuming tacit skills that may erode (A5)?
        \end{itemize}
  \item \textbf{Design.}
        For each pattern, ask how to instantiate it:
        \begin{itemize}
          \item What is the narrow, strong task we can safely
                automate (P1)?
          \item How do we express uncertainty and ``I don't know''
                states (P2)?
          \item Where are our boundaries, and how do we surface them
                in UI and logs (P3)?
          \item What are our safe horizons and re-grounding
                procedures (P4)?
          \item What does governance look like as an operator, not
                just a checklist (P5)?
\end{itemize}
\end{enumerate}